\documentclass[a4paper, 12pt]{report}
\usepackage[ngerman]{babel}
\usepackage{amsmath}  % Optional, für mathematische Formeln
\usepackage{listings}
\usepackage[table,xcdraw]{xcolor}
\usepackage{graphicx}
\usepackage[T1]{fontenc}        % Korrekte Zeichenkodierung
\usepackage{lmodern}            % Bessere Schriftqualität
\usepackage[margin=3cm]{geometry} % Ränder einstellen
\usepackage{setspace}           % Zeilenabstand
\usepackage{amsmath, amssymb}   % Mathematische Symbole
\usepackage{csquotes}           % Korrekte Zitierung mit biblatex
\usepackage[backend=biber, style=alphabetic]{biblatex} % Besseres Literaturverzeichnis
\usepackage{booktabs}           % Schönere Tabellen
\usepackage{fancyhdr}           % Kopf- und Fußzeilen
\usepackage[colorlinks=true]{hyperref} % Verlinkung in PDF

% JavaScript-Syntax-Highlighting
\lstdefinelanguage{JavaScript}{
  keywords={break, case, catch, class, const, continue, debugger, default, delete, do, else, export, extends, finally, for, function, if, import, in, instanceof, let, new, return, super, switch, this, throw, try, typeof, var, void, while, with},
  sensitive=true,
  comment=[l]{//},                     % Einzelzeilenkommentare
  comment=[s]{/*}{*/},                 % Mehrzeilige Kommentare
  string=[b]"{}",                      % String in doppelten Anführungszeichen
  morekeywords={async, await},         % Weitere Schlüsselwörter
}

\lstset{
    language=JavaScript,                  % Verwende die benutzerdefinierte JavaScript-Syntax
    backgroundcolor=\color{black!5},       % Hintergrundfarbe des Codes
    basicstyle=\ttfamily\small,            % Schriftart und -größe
    breaklines=true,                       % Zeilenumbruch bei zu langen Zeilen
    numbers=left,                          % Zeilennummern auf der linken Seite
    numberstyle=\tiny\color{gray},         % Stil der Zeilennummern
    commentstyle=\color{green}\ttfamily,    % Kommentare in grün
    keywordstyle=\color{blue}\bfseries,    % Schlüsselwörter in blau und fett
    stringstyle=\color{red},               % Strings in rot
    frame=single,                          % Rahmen um den Code
    rulecolor=\color{gray},                % Farbe des Rahmens
    breakatwhitespace=true,                % Umbrüche auch bei Leerzeichen
    tabsize=2,                             % Tabulatorgröße
    showspaces=false,                      % Keine Leerzeichensymbole
    showstringspaces=false,                % Keine Leerzeichensymbole in Strings
    morekeywords={async, await}            % Weitere Keywords für ES6 und neuer
}
\addbibresource{bibliography.bib}
\begin{document}

\title{\textit{Barrierefreiheit: Vom Gesetz zum Code} \\ 
       \textbf{Welche Auswirkungen hat das BFSG auf moderne Webanwendungen?} \\[1em]
       {Bachelor-Thesis}}

\author{Fiete Scheel} % Author
\date{\today} % Date

\maketitle

\newpage
\tableofcontents

\newpage
\chapter{Einleitung}
\input{chapters/einfuehrung.tex}

\chapter{Grundlagen}
\section{Gesetzliche Grundlagen und Standards}
\subsection{Barrierefreiheitsstärkungsgesetz}
\subsubsection{Gesetz}
Um zu gewährleisten, dass Menschen mit Behinderungen, Einschränkungen und ältere Menschen gleichermaßen an Produkten und Dienstleistungen teilhaben können, wurde das Barrierefreiheitsstärkungsgesetz \cite{bmas_bfsg} ins Leben gerufen.  

Bisher gab es für Barrierefreiheit keine einheiltlichen Regelungen, an die sich die Anbieter halten mussten. In diesem Gesetz wird die EU Richlinie zur Barrierefreiheit (European Accessibility Act, kurz: EAA) umgesetzt. Nach aktuellem Stand \textcolor{red}{müssen alle, die in der Europäischen Union etwas produzieren, verkaufen oder Dienstleistungen anbieten, nicht nur ganz unterschiedliche Anforderungen für Barrierefreiheit beachten, teilweise widersprechen sich die Anforderungen sogar. } 

Durch die klaren und einheiltichen Standards soll der Binnenmarkt gestärkt werden und eine größere Verfügbarkeit preisgünstiger barrierefreier Produkte und Dienstleistungen sicherstellen.

Am 01.03.2021 hat das Bundesministerium seinen ersten Referentenentwurf veröffentlicht. Am 24.03.2021 wurde er von der Regierung im Kabinett beschlossen und verabschiedet, ehe es am 22.07.2021 zum Abschluss kam.
Im Artikel 3 'Inkrafttreten' wird wie der Name schon sagt, auch gestgelegt, dass das Gesetz zum 28. Juni 2025 in Kraft tritt. \cite{bmas_bfsg_pdf}


\subsubsection{Verordnung}
Die Verordnung \cite{bmas_verodnung_bfsg} soll sicherstellen, dass Produkte und Dienstleistungen so gestaltet werden, dass sie für alle Menschen, insbesondere für Menschen mit Behinderungen, zugänglich und nutzbar sind. Dies umfasst sowohl technische Anforderungen als auch die Bereitstellung von Informationen und Unterstützungsdiensten. Außerdem wird auch hier noch einmal das Inkrafttreten am 28. Juni 2025 bestätigt. Die Verodnung ist am 22. Juni beschlossen worden. \cite{bmas_verordnung_bfsg_pdf}

\subsection{EU-Richtlinie}
Die EU-Richtlinie 2019/882 \cite{eu_richtlinie_2019/882}

%\chapter{Basic Elemente}
%\input{chapters/example_chapter.tex}

\printbibliography  % Hier die Bibliografie ausgeben
\end{document}
